%% 第三章

\chapter{面向低时延和低功耗的WLAN仿真平台设计与实现}
\chaptermark{面向低时延和低功耗的WLAN仿真平台设计与实现}
随着区域内通信设备数量，无线网络流量种类和规模的持续增长，未来网络对通信低时延、
能量低功耗要求尤为迫切。然而，任何预研技术在迈向落地前都需要进行前期的仿真和验证，
以评估其在实际网络环境中的性能表现，为后续的工程落地提供实验依据，可有效缩短开发流程。
因此，一个优秀且可靠的仿真平台对于网络性能评估和优化具有重要意义。本章在课题组现有的
链路级WLAN仿真平台的基础上，针对下一代WLAN的最新需求，设计和实现了粒子化接入技术、
增强型EDCA接入技术和其对应的强化学习改进模块仿真，实现了具备低时延和低功耗保障能力的下一代WLAN仿真平台。
本章首先介绍仿真平台的整体架构和设计思路，然后详细阐述低时延和低功耗保障功能的设计与实现，
最后通过仿真验证展示平台的性能表现，结果表明该仿真平台能够有效支持下一代WLAN的低时延和低功耗需求，
并且能够为后续的第四，第五章的研究提供可靠的实验基础。
\section{仿真平台架构介绍}
XXX
\section{仿真平台设计}
XXX
\subsection{整体设计}

\subsection{低时延保障功能设计}

\subsection{低功耗仿真功能设计}

\section{仿真平台实现}

\subsection{低时延保障功能实现}

\subsection{低功耗仿真功能实现}

\section{仿真平台验证}

\section{本章小结}